%% Question bank chapter file - generated by bank_tools.py build (do not edit while the bank is still being built:
%% it is regenerated from the per-test data in Question Bank\data). Every question is one bankq block:
%%   \begin{bankq}[id=..., source={...}, paper=.., qno=.., type=..., subtopic={...}, answer={...}, answer_src=key|solved,
%%                 working={...}, difficulty=1-10, diff_parts={K S I C F}, diff_note={...},
%%                 group=<passage id or empty>, also_in={duplicate copies in other tests}]
%% difficulty = min(10, K+S+I+C+F-1): K concepts 1-3, S steps 1-4, I insight 0-2, C computation 0-2, F format 0-1
%%   ...question LaTeX, options, figures (\bfig / \bfigi -> img/)...
%%   \end{bankq}
%% Passages shared by several questions are bankpassage blocks; their questions carry group=<passage id>.
\documentclass[10pt]{article}
\usepackage[a4paper,margin=1.8cm]{geometry}
\usepackage[T1]{fontenc}
\usepackage{amsmath,amssymb,mathtools}
\usepackage{graphicx}
\usepackage[export]{adjustbox}
\usepackage{xcolor,needspace,keyval}
\usepackage{enumitem}
\usepackage{booktabs,array,tabularx}
\usepackage[hidelinks]{hyperref}
\setlength{\parindent}{0pt}
\setlength{\parskip}{3pt}
\makeatletter
\define@key{bq}{id}{\def\bq@id{#1}}\define@key{bq}{source}{\def\bq@src{#1}}\define@key{bq}{paper}{\def\bq@paper{#1}}
\define@key{bq}{qno}{\def\bq@qno{#1}}\define@key{bq}{type}{\def\bq@type{#1}}\define@key{bq}{subtopic}{\def\bq@sub{#1}}
\define@key{bq}{answer}{\def\bq@ans{#1}}\define@key{bq}{answer_src}{\def\bq@asrc{#1}}\define@key{bq}{working}{\def\bq@work{#1}}
\define@key{bq}{group}{\def\bq@group{#1}}\define@key{bq}{also_in}{\def\bq@also{#1}}
\define@key{bq}{difficulty}{\def\bq@diff{#1}}\define@key{bq}{diff_parts}{\def\bq@dparts{#1}}\define@key{bq}{diff_note}{\def\bq@dnote{#1}}
\newenvironment{bankq}[1][]{\def\bq@work{}\def\bq@also{}\def\bq@group{}\def\bq@diff{?}\def\bq@dparts{}\def\bq@dnote{}\setkeys{bq}{#1}%
  \par\needspace{4\baselineskip}\medskip\noindent{\small\textbf{\bq@id}\hfill\color{gray}[\bq@src, Paper \bq@paper, Q\bq@qno]
  \bq@type\ $\cdot$ \bq@sub\ $\cdot$ \textbf{D\bq@diff}}\par\nopagebreak\relax}
  {\par{\small\color{teal}Answer: \textbf{\bq@ans}\ifx\bq@work\@empty\else\ (solved: \bq@work)\fi
  \ifx\bq@also\@empty\else\par\color{gray}Also in: \bq@also\fi}\par}
\define@key{bp}{id}{\def\bp@id{#1}}\define@key{bp}{source}{\def\bp@src{#1}}\define@key{bp}{questions}{\def\bp@qs{#1}}
\newenvironment{bankpassage}[1][]{\setkeys{bp}{#1}\par\needspace{4\baselineskip}\bigskip\noindent
  {\small\textbf{Passage \bp@id}\hfill\color{gray}[\bp@src; questions \bp@qs]}\par\nopagebreak\relax\itshape}{\par}
\makeatother
\newcommand{\bfig}[2]{\par\smallskip{\centering\includegraphics[width=#1pt,max width=0.62\linewidth,max height=0.32\textheight]{img/#2}\par}\smallskip}
\newcommand{\bfigi}[2]{\includegraphics[width=#1pt,max width=\linewidth,max height=0.32\textheight,valign=c]{img/#2}}
\newcommand{\bmissing}[1]{\textit{[figure missing in the source conversion: \detokenize{#1}]}}
\newenvironment{qoptions}{\begin{itemize}[leftmargin=1.8em,itemsep=1pt,topsep=2pt,label={}]}{\end{itemize}}
\newcommand{\opt}[2]{\item[(#1)] #2}
\providecommand{\fwstart}{}\providecommand{\fwend}{}\providecommand{\qfullsplit}{}
\newcolumntype{L}[1]{>{\raggedright\arraybackslash}p{#1}}
\newcolumntype{C}[1]{>{\centering\arraybackslash}p{#1}}
\newcolumntype{R}[1]{>{\raggedleft\arraybackslash}p{#1}}
\providecommand{\sectionhdr}[1]{}\providecommand{\subsectionhdr}[1]{}


\begin{document}
{\Large\bfseries Theory of Equations --- Question Bank}\par\smallskip
{\small Mathematics chapter M01. Questions from all converted test papers, grouped by source series. Answers: printed key where it matches the paper, otherwise solved (answer\_src=solved, with working). Exact repeats across papers appear once, with also\_in listing the other copies.}\par\medskip\hrule

\section*{Score-I papers}
\begin{bankq}[id=SI12-P2-MATH-S3-Q09, source={Score-I - Test 12 (09 Mar 2022)}, paper=2, qno=9, type=matching, subtopic={Higher-degree equations}, answer={P$\to$3, Q$\to$2, R$\to$1}, answer_src=key, difficulty=6, diff_parts={K2 S3 I1 C0 F1}, diff_note={concepts: quartic intersections, Vieta's relations | steps: h - g = 0 as a quartic; use g(0), g'(0), g''(0) for coefficients; apply the sum condition; solve a, c; sum of roots | I: sum i/x\_i condition}]
Let $h(x)=4x^4-ax^3$ and $g(x)$ is a quadratic polynomial satisfying $g(0)+6=g'(0)-c=g''(0)+2b=0$ (where $a,b,c\in R$). If $y=h(x)$ and $y=g(x)$ intersect at 4 distinct points with non negative abscissae $x_i$, $i=1,2,3,4$ such that $\displaystyle\sum_{i=1}^4\dfrac i{x_i}=8$.
\begin{center}\footnotesize
\begin{tabular}{|p{0.6cm}|p{1.6cm}|p{0.6cm}|p{1.0cm}|}
\hline
 & \centering List-I & & \centering List-II\tabularnewline
\hline
(P) & $a$ & (1) & 5\\
(Q) & $c$ & (2) & 25\\
(R) & $\displaystyle\sum_{i=1}^4x_i$ & (3) & 20\\
\hline
\end{tabular}
\end{center}
\end{bankq}

\section*{TOAS papers}
\begin{bankq}[id=TOAS01-P1-MATH-S1-Q02, source={TOAS - Test 01 (30 Jan 2022)}, paper=1, qno=2, type=single, subtopic={Location of roots}, answer={B}, answer_src=key, difficulty=1, diff_parts={K1 S1 I0 C0 F0}, diff_note={concepts: position of a number relative to the roots | steps: f(p) < 0 with p = 8 - a; solve for a and read the maximum integer p}]
If product of roots `$p$' of equation $x^2-2ax+8-a=0$ lies between its roots, then maximum integral value of `$p$' is -
\begin{qoptions}
\opt{A}{1}
\opt{B}{4}
\opt{C}{6}
\opt{D}{7}
\end{qoptions}
\end{bankq}
\begin{bankq}[id=TOAS01-P1-MATH-S2-Q06, source={TOAS - Test 01 (30 Jan 2022)}, paper=1, qno=6, type=multi, subtopic={Quadratics}, answer={A,C}, answer_src=key, difficulty=2, diff_parts={K1 S1 I0 C0 F1}, diff_note={concepts: roots of a quadratic | steps:  | alpha - beta}]
If $f(x)=2x^2+kx+50$, $k\ge20$ and magnitude of difference of the roots of the equation $f(x)=0$ is 8, then which of the following can be true -
\begin{qoptions}
\opt{A}{minimum value of $f(x)$ is $-32$}
\opt{B}{minimum value of $f(x)$ is $-64$}
\opt{C}{graph of $f(x)$ cuts positive y-axis}
\opt{D}{graph of $f(x)$ cuts negative y-axis}
\end{qoptions}
\end{bankq}
\begin{bankq}[id=TOAS01-P1-MATH-S4-Q08, source={TOAS - Test 01 (30 Jan 2022)}, paper=1, qno=8, type=numerical, subtopic={Higher-degree equations}, answer={729}, answer_src=key, difficulty=3, diff_parts={K1 S1 I1 C0 F1}, diff_note={concepts: polynomial roots | steps: product of (alpha - 2) equals -f(2) = 9; cube it | I: evaluate the polynomial at 2}]
If $\alpha,\beta,\gamma$ are the roots of the equation $x^3-5x^2+3=0$, then the value of $(\alpha-2)^3(\beta-2)^3(\gamma-2)^3$ is
\end{bankq}
\begin{bankq}[id=TOAS01-P2-MATH-S1-Q03, source={TOAS - Test 01 (30 Jan 2022)}, paper=2, qno=3, type=single, subtopic={Quadratics}, answer={D}, answer_src=key, difficulty=1, diff_parts={K1 S1 I0 C0 F0}, diff_note={concepts: inequality with modulus | steps: put t =  | x - 2}]
Sum of all possible integral solutions of the inequality $(x-2)^2-5|x-2|+6<0$ is -
\begin{qoptions}
\opt{A}{4}
\opt{B}{5}
\opt{C}{7}
\opt{D}{8}
\end{qoptions}
\end{bankq}
\begin{bankq}[id=TOAS01-P2-MATH-S1-Q14, source={TOAS - Test 01 (30 Jan 2022)}, paper=2, qno=14, type=single, subtopic={Quadratics}, answer={A}, answer_src=key, difficulty=1, diff_parts={K1 S1 I0 C0 F0}, diff_note={concepts: identity as a quadratic | steps: all three coefficients zero gives p = 1}]
If the equation $(p^2-3p+2)x^2+(p-1)^2x+(p^2-4p+3)=0$ is satisfied for more than two different values of $x$, then $p$ is equal to -
\begin{qoptions}
\opt{A}{1}
\opt{B}{2}
\opt{C}{1,2,3}
\opt{D}{2,3}
\end{qoptions}
\end{bankq}
\begin{bankq}[id=TOAS01-P2-MATH-S1-Q15, source={TOAS - Test 01 (30 Jan 2022)}, paper=2, qno=15, type=single, subtopic={Quadratics}, answer={C}, answer_src=key, difficulty=2, diff_parts={K1 S1 I1 C0 F0}, diff_note={concepts: roots of a quadratic | steps: alpha beta = 64/27; expression equals (alpha + beta)/(alpha beta)\^{}(1/3) | I: combine into one fraction}]
If $\alpha$ and $\beta$ are the roots of equation $27x^2-4x+64=0$, then the value of $\left(\dfrac{\alpha^{2/3}}{\beta^{1/3}}+\dfrac{\beta^{2/3}}{\alpha^{1/3}}\right)$ is -
\begin{qoptions}
\opt{A}{$\dfrac13$}
\opt{B}{$\dfrac14$}
\opt{C}{$\dfrac19$}
\opt{D}{$\dfrac1{16}$}
\end{qoptions}
\end{bankq}
\begin{bankq}[id=TOAS01-P2-MATH-S2-Q09, source={TOAS - Test 01 (30 Jan 2022)}, paper=2, qno=9, type=numerical, subtopic={Higher-degree equations}, answer={5}, answer_src=solved, working={$b+c=-a$, $a^3-2=a-1$; $\frac{a^2}{a-1}=a+1+\frac1{a-1}$; $\sum(a+1)=3$, $\sum\frac1{a-1}=-\frac{p^{\prime}(1)}{p(1)}=-\frac{2}{-1}=2$; total 5}, difficulty=5, diff_parts={K1 S3 I1 C0 F1}, diff_note={concepts: symmetric functions of roots | steps: b + c = -a; a\^{}3 - 2 = a - 1; divide a\^{}2/(a-1); sum of (a + 1); sum of 1/(a-1) from the derivative | I: partial fraction form}]
If $a,b,c$ are the roots of the equation $x^3-x-1=0$, then the value of $\dfrac{(b+c)^2}{a^3-2}+\dfrac{(c+a)^2}{b^3-2}+\dfrac{(a+b)^2}{c^3-2}$ is
\end{bankq}
\begin{bankq}[id=TOAS02-P1-MATH-S2-Q11, source={TOAS - Test 02 (13 Feb 2022)}, paper=1, qno=11, type=multi, subtopic={Higher-degree equations}, answer={A,C,D}, answer_src=key, difficulty=6, diff_parts={K2 S2 I2 C0 F1}, diff_note={concepts: cubic equation, inverse-function symmetry | steps: recognise (x\^{}3 + 3/8)\^{}3; set t = x\^{}3 + 3/8 so t\^{}3 = x - 3/8; t and x satisfy inverse relations so t = x; solve x\^{}3 - x + 3/8 = 0 | I: inverse function trick}]
The values of $x$ satisfying the equation $x^9+\dfrac98x^6+\dfrac{27}{64}x^3-x+\dfrac{219}{512}=0$ is -
\begin{qoptions}
\opt{A}{$\dfrac{-1+\sqrt{13}}{4}$}
\opt{B}{$-\dfrac12$}
\opt{C}{$\dfrac{-1-\sqrt{13}}{4}$}
\opt{D}{$\dfrac12$}
\end{qoptions}
\end{bankq}
\begin{bankq}[id=TOAS05-P1-MATH-S1-Q08, source={TOAS - Test 05 (03 Apr 2022)}, paper=1, qno=8, type=multi, subtopic={Higher-degree equations}, answer={A,B,C,D}, answer_src=key, difficulty=5, diff_parts={K1 S2 I2 C0 F1}, diff_note={concepts: polynomial root estimate | steps: group as (x - 1)(x\^{}10 + x\^{}7 + x\^{}4 + x) = 20; factor x\^{}9 + x\^{}6 + x\^{}3 + 1; relate alpha\^{}12 to a bound; compare with the options | I: factor the polynomial}]
If $\alpha$ ($\alpha>0$) is a root of $x^{11}-x^{10}+x^8-x^7+x^5-x^4+x^2-x-20=0$ then $\alpha^{12}$ is greater than
\begin{qoptions}
\opt{A}{59}
\opt{B}{61}
\opt{C}{58}
\opt{D}{57}
\end{qoptions}
\end{bankq}
\begin{bankq}[id=TOAS05-P1-MATH-S2-Q09, source={TOAS - Test 05 (03 Apr 2022)}, paper=1, qno=9, type=single, subtopic={Higher-degree equations}, answer={D}, answer_src=key, difficulty=3, diff_parts={K1 S2 I0 C1 F0}, diff_note={concepts: roots of higher-degree equations | steps: solve the equation in List-I item I; item II; item III; item IV}]
Match List-I with List-II and select the correct answer using the code given below the list.
{\small\setlength{\tabcolsep}{3pt}\renewcommand{\arraystretch}{1.25}\begin{center}\begin{tabular}{|L{0.8cm}|L{4.2cm}|L{0.65cm}|L{0.75cm}|}
\hline
\multicolumn{2}{|c|}{\parbox{4.6cm}{\centering List-I\\Equations}} & \multicolumn{2}{c|}{\parbox{1.5cm}{\centering\vphantom{\rule{0pt}{2ex}}List-II\\No. of solutions\vphantom{\rule[-1ex]{0pt}{1ex}}}}\\
\hline
(I) & $x^2\tan x=1$\ \ $x\in\left(\dfrac\pi2,\dfrac{3\pi}2\right)$ & (P) & 5\\
\hline
(II) & $2^{\cos x}=|\sin x|$, $x\in[0,2\pi]$ & (Q) & 1\\
\hline
(III) & If $f(x)$ is a polynomial of degree 5 with real coefficients such that $f(|x|)=0$ has 8 real roots then, roots of $f(x)=0$ & (R) & 3\\
\hline
(IV) & $7^{|x|}\left(|5-|x||\right)=1$ & (S) & 4\\
\hline
\multicolumn{2}{|c|}{} & (T) & 0\\
\hline
\multicolumn{2}{|c|}{} & (U) & 10\\
\hline
\end{tabular}\end{center}}
Which of the following is the only \textbf{CORRECT} combination ?
\begin{qoptions}
\opt{A}{I $\to$ P}
\opt{B}{II $\to$ Q}
\opt{C}{III $\to$ R}
\opt{D}{IV $\to$ S}
\end{qoptions}
\end{bankq}
\begin{bankq}[id=TOAS05-P1-MATH-S2-Q10, source={TOAS - Test 05 (03 Apr 2022)}, paper=1, qno=10, type=single, subtopic={Higher-degree equations}, answer={A}, answer_src=key, difficulty=3, diff_parts={K1 S2 I0 C1 F0}, diff_note={concepts: roots of higher-degree equations | steps: solve the equation in List-I item I; item II; item III; item IV}]
Match List-I with List-II and select the correct answer using the code given below the list.
{\small\setlength{\tabcolsep}{3pt}\renewcommand{\arraystretch}{1.25}\begin{center}\begin{tabular}{|L{0.8cm}|L{4.2cm}|L{0.65cm}|L{0.75cm}|}
\hline
\multicolumn{2}{|c|}{\parbox{4.6cm}{\centering List-I\\Equations}} & \multicolumn{2}{c|}{\parbox{1.5cm}{\centering\vphantom{\rule{0pt}{2ex}}List-II\\No. of solutions\vphantom{\rule[-1ex]{0pt}{1ex}}}}\\
\hline
(I) & $x^2\tan x=1$\ \ $x\in\left(\dfrac\pi2,\dfrac{3\pi}2\right)$ & (P) & 5\\
\hline
(II) & $2^{\cos x}=|\sin x|$, $x\in[0,2\pi]$ & (Q) & 1\\
\hline
(III) & If $f(x)$ is a polynomial of degree 5 with real coefficients such that $f(|x|)=0$ has 8 real roots then, roots of $f(x)=0$ & (R) & 3\\
\hline
(IV) & $7^{|x|}\left(|5-|x||\right)=1$ & (S) & 4\\
\hline
\multicolumn{2}{|c|}{} & (T) & 0\\
\hline
\multicolumn{2}{|c|}{} & (U) & 10\\
\hline
\end{tabular}\end{center}}
Which of the following is the only \textbf{INCORRECT} combination ?
\begin{qoptions}
\opt{A}{I $\to$ R}
\opt{B}{II $\to$ S}
\opt{C}{III $\to$ P}
\opt{D}{IV $\to$ Q}
\end{qoptions}
\end{bankq}
\begin{bankq}[id=TOAS06-P1-MATH-S3-Q02, source={TOAS - Test 06 (17 Apr 2022)}, paper=1, qno=2, type=numerical, subtopic={Location of roots}, answer={6.00}, answer_src=key, difficulty=4, diff_parts={K1 S2 I1 C0 F1}, diff_note={concepts: zeros of an even polynomial | steps: P(0) = -1, P(1) = 1, P(2) = -1; sign changes in x > 0; evenness doubles | I: use evenness}]
Let $P(x)=x^{10}+ax^8+bx^6+cx^4+dx^2-1$ be polynomial with real coefficients. If $P(1)=1$, $P(2)=-1$ then the minimum number of real zeroes of $P(x)$ is
\end{bankq}
\begin{bankq}[id=TOAS08-P2-MATH-S3-Q02, source={TOAS - Test 08 (15 May 2022)}, paper=2, qno=2, type=numerical, subtopic={Higher-degree equations}, answer={3.20}, answer_src=key, difficulty=3, diff_parts={K1 S1 I1 C0 F1}, diff_note={concepts: real root of a symmetric sum of odd powers | steps: the sum is increasing and odd about the mean value 1010 | I: symmetry about the mean}]
If the real root of the equation $(x-1)^{2019}+(x-2)^{2019}+(x-3)^{2019}+\ldots+(x-2019)^{2019}=0$ is $\lambda$, then $\dfrac{\lambda}{315}$ equals
\end{bankq}
\begin{bankq}[id=TOAS09-P1-MATH-S2-Q07, source={TOAS - Test 09 (30 May 2022)}, paper=1, qno=7, type=multi, subtopic={Higher-degree equations}, answer={D}, answer_src=key, difficulty=6, diff_parts={K2 S2 I2 C0 F1}, diff_note={concepts: roots of a cubic, double limit | steps: roots of x\^{}3 - 3x + 1 = 0 are all real and irrational; cos(n! pi x) for irrational x; limit in m gives 0 | I: irrational numbers}]
Let $x_1,x_2,x_3$ be the roots of the equation $x(3-x^2)=1$, then
\begin{qoptions}
\opt{A}{exactly one of $x_1,x_2,x_3$ is real}
\opt{B}{exactly one of $x_1,x_2,x_3$ is rational}
\opt{C}{the value of $\displaystyle\sum_{i=1}^{3}\left(\lim_{m\to\infty}\left(\lim_{n\to\infty}\cos(\pi n!\,x_i)\right)^{m}\right)$ where $m,n\in N$, is 3}
\opt{D}{the value of $\displaystyle\sum_{i=1}^{3}\left(\lim_{m\to\infty}\left(\lim_{n\to\infty}\cos(\pi n!\,x_i)\right)^{m}\right)$ where $m,n\in N$, is 0}
\end{qoptions}
\end{bankq}
\begin{bankq}[id=TOAS09-P1-MATH-S4-Q08, source={TOAS - Test 09 (30 May 2022)}, paper=1, qno=8, type=numerical, subtopic={Quadratics}, answer={2}, answer_src=key, difficulty=2, diff_parts={K1 S1 I0 C0 F1}, diff_note={concepts: identity in n fixes constants | steps: coefficient of n and the constant term give two linear equations; solve}]
For $\ell,m,n\in R$, $p$ and $q$ are the value of $\ell$ and $m$ respectively such that equation $\ell(13+11n)+m(19+17n)-4(8+7n)=0$ is satisfied $\forall\ n\in R$ then $(p+q)$ equals
\end{bankq}
\begin{bankq}[id=TOAS09-P2-MATH-S1-Q18, source={TOAS - Test 09 (30 May 2022)}, paper=2, qno=18, type=single, subtopic={Higher-degree equations}, answer={A}, answer_src=key, difficulty=1, diff_parts={K1 S1 I0 C0 F0}, diff_note={concepts: double root of a cubic | steps: f'(x) = 0 at x = -1 and 3; c from f = 0 there}]
The expression $x^3-3x^2-9x+c$ can be written in the form $(x-a)^2(x-b)$ if the values of c is
\begin{qoptions}
\opt{A}{$-5$ or 27}
\opt{B}{5 or $-27$}
\opt{C}{5 or 27}
\opt{D}{$-5$ or $-27$}
\end{qoptions}
\end{bankq}
\begin{bankq}[id=TOAS10-P1-MATH-S1-Q07, source={TOAS - Test 10 (12 Jun 2022)}, paper=1, qno=7, type=multi, subtopic={Higher-degree equations}, answer={A,C}, answer_src=key, difficulty=5, diff_parts={K2 S2 I1 C0 F1}, diff_note={concepts: symmetric sums, a cubic with real roots | steps: reduce to ab + bc + ca = 5; (a+b)(b+c)(c+a) = 20 - abc; range of abc from the cubic's extrema | I: three real roots bound abc}]
If $a+b+c=4$ and $a^2+b^2+c^2+3(ab+bc+ac)=21$, where $a,b,c\in\mathbb{R}$, then
\begin{qoptions}
\opt{A}{minimum value of $(a+b)(b+c)(a+c)$ is 18}
\opt{B}{minimum value of $(a+b)(b+c)(a+c)$ is 24}
\opt{C}{maximum value of $(a+b)(b+c)(a+c)$ is $\dfrac{490}{27}$}
\opt{D}{maximum value of $(a+b)(b+c)(a+c)$ is $\dfrac{471}{27}$}
\end{qoptions}
\end{bankq}
\begin{bankq}[id=TOAS12-P1-MATH-S1-Q08, source={TOAS - Test 12 (10 Jul 2022)}, paper=1, qno=8, type=multi, subtopic={Location of roots}, answer={D}, answer_src=key, difficulty=4, diff_parts={K1 S2 I1 C0 F1}, diff_note={concepts: zeros of a function using parity | steps: a x\^{}9 + b sin x + c x\^{}2 sgn x is odd; the tanh term is odd; f is odd so zeros are symmetric; IVT between the given values | I: f is an odd function}]
Let $f(x)=ax^9+b\sin x+cx^2\operatorname{sgn}(x)+\dfrac{e^x-e^{-x}}{e^x+e^{-x}}$ be defined on set of real number (where $a>0$, $b,c\in\mathbb{R}$). If $f(-5)=5$, $f(-2)=-3$, then the minimum number of zeroes of the equation $f(x)=0$, is
\begin{qoptions}
\opt{A}{1}
\opt{B}{2}
\opt{C}{3}
\opt{D}{5}
\end{qoptions}
\end{bankq}
\begin{bankq}[id=TOAS13-P1-MATH-S3-Q04, source={TOAS - Test 13 (24 Jul 2022)}, paper=1, qno=4, type=numerical, subtopic={Higher-degree equations}, answer={8}, answer_src=solved, working={Degree $n$: $2^n=16$, so $n=4$. Putting $x=16,1,2,4$ gives $f(16)=f(2)=f(4)=f(8)=0$, so $f=(x-2)(x-4)(x-8)(x-16)$, $f(0)=1024$, $f(0)/128=8$}, difficulty=4, diff_parts={K1 S2 I1 C0 F1}, diff_note={concepts: polynomial functional equation | steps: degree from leading coefficients is 4; substitute x = 16, 1, 2, 4 to get the roots; f(0)/128}]
Let $f(x)$ be a polynomial with leading coefficient unity and satisfying the equation $(x-16)f(2x)=16(x-1)f(x)$ then $\dfrac{f(0)}{128}$ is equal to \underline{\hspace{1.5cm}}.
\end{bankq}
\begin{bankq}[id=TOAS13-P1-MATH-S3-Q06, source={TOAS - Test 13 (24 Jul 2022)}, paper=1, qno=6, type=numerical, subtopic={Higher-degree equations}, answer={0}, answer_src=solved, working={$\dfrac{d}{dx}\dfrac{f\prime}{f}=\dfrac{ff\prime\prime-f\prime^2}{f^2}=-\sum\dfrac1{(x-r_i)^2}<0$, so it never vanishes off the roots; at a root the expression is $-f\prime^2\ne0$ (roots are simple). No real roots}, difficulty=4, diff_parts={K1 S1 I2 C0 F1}, diff_note={concepts: logarithmic derivative of a polynomial | steps: derivative of f'/f is minus the sum of squares so it never vanishes; at a root the expression is minus f'\^{}2}]
$f(x)$ is a fifth order polynomial in $x$ with every root of $f(x)=0$ is real and distinct, then the number of real roots of $f''(x)\,f(x)-(f'(x))^2=0$ is equal to \underline{\hspace{1.5cm}}.
\end{bankq}
\begin{bankq}[id=TOAS15-P1-MATH-S2-Q07, source={TOAS - Test 15 (22 Aug 2022)}, paper=1, qno=7, type=multi, subtopic={Higher-degree equations}, answer={A,C,D}, answer_src=key, difficulty=6, diff_parts={K2 S2 I2 C0 F1}, diff_note={concepts: symmetric system from a cube root equation, cubic factorisation | steps: set y as the cube root; get the symmetric pair x\^{}3 + 1 = 2y and y\^{}3 + 1 = 2x; x = y gives the cubic; read a b c d}]
If the solutions of the equation $x^3+1=2\sqrt[3]{2x-1}$ are $\left\{a,\dfrac{b\pm\sqrt c}{d}\right\}$, $(a,\ b\in I\ \&\ c,d$ are prime) then which of the following is(are) correct ?
\begin{qoptions}
\opt{A}{$a+b^3+c-d=3$}
\opt{B}{$a+b=1$}
\opt{C}{$c+d=7$}
\opt{D}{$c-d=3$}
\end{qoptions}
\end{bankq}
\begin{bankq}[id=TOAS19-P1-MATH-S2-Q09, source={TOAS - Test 19 (05 Feb 2023)}, paper=1, qno=9, type=single, subtopic={Higher-degree equations}, answer={B}, answer_src=key, difficulty=4, diff_parts={K1 S2 I1 C1 F0}, diff_note={concepts: symmetric functions of the roots | steps: sum of squares; value at 1 and -1; reduction of the cube expression; sum of reciprocals}]
Answer the following by appropriately matching the lists based on the information given in the paragraph

$\alpha_1,\alpha_2,\alpha_3,\alpha_4$ and $\alpha_5$ are all roots of $x^5-4x^3-3x^2+2x-2=0$.

There are some expression given in List-I whose values are given in List-II below :
\begin{center}\footnotesize\setlength{\tabcolsep}{2pt}
\begin{tabular}{|L{0.8cm}L{4.4cm}|L{0.6cm}L{0.6cm}|}\hline
\multicolumn{2}{|c|}{List-I}&\multicolumn{2}{c|}{List-II}\\\hline
(I)&$\alpha_1^2+\alpha_2^2+\alpha_3^2+\alpha_4^2+\alpha_5^2$&(P)&1\\\hline
(II)&$\left|(1-\alpha_1^2)(1-\alpha_2^2)\dots(1-\alpha_5^2)\right|$&(Q)&8\\\hline
(III)&$\bigl|(\alpha_1^3-4\alpha_1-3)(\alpha_2^3-4\alpha_2-3)$ $\dots(\alpha_5^3-4\alpha_5-3)\bigr|$&(R)&12\\\hline
(IV)&$\dfrac1{\alpha_1}+\dfrac1{\alpha_2}+\dots\dfrac1{\alpha_5}$&(S)&24\\\hline
&&(T)&36\\\hline
&&(U)&48\\\hline
\end{tabular}\end{center}
Which of the following is the only correct combination ?
\begin{qoptions}
\opt{A}{I $\to$ R}
\opt{B}{II $\to$ S}
\opt{C}{III $\to$ T}
\opt{D}{IV $\to$ Q}
\end{qoptions}
\end{bankq}
\begin{bankq}[id=TOAS19-P1-MATH-S2-Q10, source={TOAS - Test 19 (05 Feb 2023)}, paper=1, qno=10, type=single, subtopic={Higher-degree equations}, answer={C}, answer_src=key, difficulty=4, diff_parts={K1 S2 I1 C1 F0}, diff_note={concepts: symmetric functions of the roots | steps: sum of squares; value at 1 and -1; reduction of the cube expression; sum of reciprocals}]
Answer the following by appropriately matching the lists based on the information given in the paragraph

$\alpha_1,\alpha_2,\alpha_3,\alpha_4$ and $\alpha_5$ are all roots of $x^5-4x^3-3x^2+2x-2=0$.

There are some expression given in List-I whose values are given in List-II below :
\begin{center}\footnotesize\setlength{\tabcolsep}{2pt}
\begin{tabular}{|L{0.8cm}L{4.4cm}|L{0.6cm}L{0.6cm}|}\hline
\multicolumn{2}{|c|}{List-I}&\multicolumn{2}{c|}{List-II}\\\hline
(I)&$\alpha_1^2+\alpha_2^2+\alpha_3^2+\alpha_4^2+\alpha_5^2$&(P)&1\\\hline
(II)&$\left|(1-\alpha_1^2)(1-\alpha_2^2)\dots(1-\alpha_5^2)\right|$&(Q)&8\\\hline
(III)&$\bigl|(\alpha_1^3-4\alpha_1-3)(\alpha_2^3-4\alpha_2-3)$ $\dots(\alpha_5^3-4\alpha_5-3)\bigr|$&(R)&12\\\hline
(IV)&$\dfrac1{\alpha_1}+\dfrac1{\alpha_2}+\dots\dfrac1{\alpha_5}$&(S)&24\\\hline
&&(T)&36\\\hline
&&(U)&48\\\hline
\end{tabular}\end{center}
Which of the following is the only correct combination ?
\begin{qoptions}
\opt{A}{I $\to$ R}
\opt{B}{II $\to$ T}
\opt{C}{III $\to$ U}
\opt{D}{IV $\to$ R}
\end{qoptions}
\end{bankq}
\end{document}
